\chapter{Gear, Magic Items, and Crafting}
\label{ch:gear}

\section{The Philosophy of Gear}
In Fate’s Edge, your character’s possessions are described in the fiction, not tracked on a ledger. You have enough coin, barter, or credit to afford:
\begin{itemize}
  \item Ordinary meals, lodging, and travel rations
  \item Basic tools (rope, lantern, writing kit, lockpicks)
  \item Small bribes (a drink, a gift)
  \item Routine transportation (ferry, caravan passage)
\end{itemize}
If the GM ever asks “Do you have enough?” the answer is usually **yes** – unless scarcity is the point of the scene.

\section{Mundane Equipment}
\label{sec:mundane-gear}

\subsection{Weapons and Armor}
Weapons and armor are covered in \S\ref{sec:weapons-armor} (Core Mechanics). Briefly:
\begin{itemize}
  \item \textbf{Light weapons (4 XP):} fast, concealable.
  \item \textbf{Medium weapons (8 XP):} balanced.
  \item \textbf{Heavy weapons (12 XP):} slow, powerful.
  \item \textbf{Light armor (4 XP):} converts 1 Harm → 1 Fatigue.
  \item \textbf{Medium armor (8 XP):} converts 2 Harm → 1 Fatigue, –1 die physical skills.
  \item \textbf{Heavy armor (12 XP):} converts 3 Harm → 2 Fatigue, –2 dice physical skills, no sprint in rough terrain.
\end{itemize}
You may acquire these with XP at character creation or later.

\subsection{Adventuring Gear}
You are assumed to have the following unless you are in a survival situation:
\begin{itemize}
  \item Backpack, waterskin, 1d6 days of rations, a knife, flint and steel, a small lantern or candle
  \item Tools for any skill you have at level 1+ (e.g., lockpicks for Subterfuge, healer’s kit for Medicine)
  \item A set of clothing appropriate to your culture
\end{itemize}
If you need something unusual (a climber’s kit, a disguise, a rare herb), you may:
\begin{itemize}
  \item Spend 1 Boon to declare you already have it (narrative permission).
  \item Spend 1 XP to acquire it (as a one‑time resource).
  \item Acquire it as a minor asset (4 XP) if you want permanent access.
\end{itemize}

\section{Magic Items}
Magic items are not bought with gold. They are acquired by spending XP, completing quests, or crafting. A magic item is a **permanent asset** that grants a mechanical benefit. It requires **upkeep** (see below).

\subsection{Cost Basis: Talent Equivalents}
A magic item’s effect is roughly equivalent to a **Talent** (see \ref{ch:talents}). Use the Talent XP cost as the base price, then apply a **discount** because items can be lost, stolen, or broken.

\begin{center}
\begin{tabular}{lcc}
\toprule
\textbf{Effect Power} & \textbf{Talent Cost} & \textbf{Magic Item Cost (XP)} \\
\midrule
Minor (e.g., +1 die to a narrow skill once/scene) & 2–3 XP & 2 XP \\
Major (e.g., +1 die to a broad skill always) & 4–6 XP & 4 XP \\
Prestige (e.g., once/session reroll a failed save) & 7–10 XP & 6 XP \\
Epic (e.g., teleport once/day) & 11+ XP & 8 XP \\
\bottomrule
\end{tabular}
\end{center}
\emph{Discount reason:} A Talent is permanently attached to your character; an item can be taken away. The discount also covers the need for upkeep.

\subsection{Sample Magic Items}

\begin{longtable}{p{3.5cm}p{3.5cm}p{7cm}}
\toprule
\textbf{Name} & \textbf{Cost (XP)} & \textbf{Effect} \\
\midrule
Emberfang Dagger & 2 & When drawn, glows faintly if undead are within Near. Once per scene, gain +1 die to detect hidden spirits. \\
Storm Lantern & 4 & Once per session, as an action, create a gust of wind that extinguishes all mundane flames in Near and gives +1 Position to disengage. \\
Cloak of the Unseen & 4 & +1 die to Stealth in dim light. Once per scene, reroll a failed Stealth check. \\
Vial of the Last Breath & 2 (consumable) & Heal 1 Fatigue when drunk. One use. \\
Pendant of the Witness & 6 & Once per session, when you detect a lie, the liar suffers 1 Fatigue (no save). \\
Raven’s Key & 8 & Once per day, as an action, unlock any mundane lock or open any door that is not magically sealed. \\
\bottomrule
\end{longtable}

\subsection{Attunement and Upkeep}
A character may be \textbf{attuned} to up to 3 magic items at a time (attunement takes a short rest). During each downtime, you must pay **upkeep** for each magic item you keep attuned:
\begin{itemize}
  \item \textbf{Efficient:} Pay XP equal to the item’s cost divided by 3 (rounded up, minimum 1). You spend minimal effort.
  \item \textbf{Intensive:} Pay 1 XP and describe a downtime scene where you maintain the item (cleaning, recharging, appeasing its spirit).
\end{itemize}
If you neglect upkeep, the item becomes **Neglected** (no benefit until you pay the upkeep retroactively). If neglected for two consecutive downtimes, it becomes **Compromised** (requires a quest to restore).

\subsection{Losing Magic Items}
If you lose a magic item (theft, destruction, sacrifice), you lose the XP you spent on it. The GM may allow you to recover it as a story hook. You may also choose to “retire” an item (sell it, trade it, give it away) and regain half its XP cost (rounded down) to spend on something else.

\section{Crafting Magic Items}
\label{sec:crafting}

Crafting a permanent magic item is a **campaign‑level project**. It requires:
\begin{itemize}
  \item A character with relevant skills: \textbf{Arcana 3+} and \textbf{Craft 2+} (or appropriate specialty).
  \item A workspace (e.g., a workshop asset, a consecrated forge).
  \item Rare components (obtained through adventuring – the GM will decide).
  \item A **Downtime Project Clock** (6–12 segments). Each segment requires a scene or a downtime action.
  \item The XP cost of the item (see table above) – paid when the clock is completed.
\end{itemize}

\subsection{Crafting Procedure}
\label{subsec:crafting-procedure}
\begin{enumerate}
  \item \textbf{Design:} Describe what the item does. The GM sets the Tier (Minor/Major/Prestige/Epic) and XP cost.
  \item \textbf{Quest:} The GM frames a short arc to obtain a rare component (e.g., “a phoenix feather”, “the heart of a storm elemental”).
  \item \textbf{Downtime Clock:} Roll \textbf{Wits + Craft} or \textbf{Wits + Arcana} each downtime. Each success advances the clock by 1. Partial successes create complications (lost components, attracting attention). On a failure, the clock may tick backwards.
  \item \textbf{Completion:} When the clock fills, you pay the XP and the item is complete. The GM may add a final, optional cost (e.g., “it also consumes a memory”).
\end{enumerate}

\subsection{Crafting Consumables}
One‑use items (potions, scrolls, runestones) are much easier. Crafting a consumable:
\begin{itemize}
  \item Takes a single downtime action (no clock).
  \item Costs \textbf{1 XP} per consumable (or 2 XP for a powerful one).
  \item Requires an appropriate skill roll (DV 3) – on a failure, the components are wasted.
  \item No upkeep required.
\end{itemize}
Examples: \emph{Healing potion} (restores 1 Fatigue), \emph{Sight scroll} (grants +1 die to Investigation for one scene), \emph{Thunderstone} (ranged attack, Harm 1 to a small area).

\section{Upgrading Mundane Gear}
You may also spend XP to give a mundane item a permanent magical property. Use the same costs as magic items. The process is faster than crafting from scratch (no rare component quest, just a downtime clock of 4–6 segments).

\section{Summary for Players}
\begin{itemize}
  \item Ordinary gear is free – describe what you carry.
  \item Weapons and armor cost XP (light 4, medium 8, heavy 12).
  \item Magic items cost XP based on Talent equivalents (2–8 XP).
  \item Magic items require upkeep (XP or downtime scene) each downtime.
  \item Crafting a magic item is a long project involving a clock, rare components, and XP.
\end{itemize}

\begin{tcolorbox}[colback=green!5!white,colframe=green!70!black,title=The Gray Wanderer’s Tip]
“A magic sword is not a toy. It is a creditor. It demands maintenance, attention, and sometimes a blood‑price. If you cannot afford the upkeep, do not pick it up. The edge is not worth the interest.”
\end{tcolorbox}
